\section{Vorläufige Gliederung}\label{gliederung}

\subsection{1 Einleitung}
1.1 Motivation und Relevanz
1.2 Problemstellung
1.3 Zielsetzung und Forschungsfragen
1.4 Struktur der Arbeit

\subsection{2 Theoretische Grundlagen}
2.1 IoT-Sicherheit: Architektur, Schwachstellen und Bedrohungslandschaft
2.2 Klassische versus ML-basierte Intrusion Detection Systeme
2.3 Maschinelles Lernen für IoT-IDS: Supervised Learning und Ensemble-Methoden
2.4 Klassenungleichgewicht in ML-Anwendungen: Ursachen, Auswirkungen und Gegenmaßnahmen
2.5 Dimensionsreduktion mit PCA: Grundlagen und Einsatz in IDS
2.6 Erklärbare Künstliche Intelligenz (XAI): SHAP als post-hoc-Erklärungsverfahren

\subsection{3 Verwandte Arbeiten}
3.1 ML-basierte Angriffserkennung auf dem CICIoT2023-Datensatz
3.2 XAI-Ansätze in Intrusion Detection Systemen
3.3 Behandlung von Klassenungleichgewicht in IDS-Studien
3.4 Abgrenzung der eigenen Arbeit vom Stand der Forschung

\subsection{4 Datensatz und Vorverarbeitung}
4.1 Beschreibung des CICIoT2023-Datensatzes
4.2 Angriffsklassen und Klassenstruktur
4.3 Klassenungleichgewicht: Analyse und Auswirkungen
4.4 Vorverarbeitungspipeline: Normalisierung und Klassenkondensierung
4.5 PCA-Dimensionsreduktion: Vorgehen und Parameterwahl

\subsection{5 Methodik}
5.1 Zweistufiger hierarchischer Klassifikationsansatz
5.2 Gradient-Boosting-Modell: Architektur und Hyperparameteroptimierung
5.3 Evaluationsmetriken: Balanced Accuracy, F1-Score, Konfusionsmatrix
5.4 Ablation Study: PCA versus vollständiger Merkmalsraum
5.5 SHAP-Integration: globale und lokale Erklärungsebene

\subsection{6 Experimentelle Evaluation}
6.1 Experimentelles Setup und Reproduzierbarkeit
6.2 Ergebnisse der binären Klassifikation (Stufe 1)
6.3 Ergebnisse der Mehrklassenklassifikation (Stufe 2)
6.4 Ergebnisse der Ablation Study
6.5 SHAP-Analyse: Feature-Importance pro Angriffsklasse
6.6 Semantische Validierung der SHAP-Erklärungen

\subsection{7 Diskussion}
7.1 Interpretation der Ergebnisse im Licht der Forschungsfragen
7.2 Vergleich mit Raturi et al. (2026) und Almahaqeri et al. (2026)
7.3 Bedeutung der Balanced Accuracy für die IDS-Bewertung
7.4 Grenzen der Arbeit

\subsection{8 Fazit und Ausblick}
8.1 Zusammenfassung der Beiträge
8.2 Beantwortung der Forschungsfragen
8.3 Offene Fragen und zukünftige Forschungsrichtungen
